% tags: Tag One, Tag Two
% summary: One or two sentences describing the note — shown on the notes index.
\documentclass[11pt]{article}
\usepackage{amsmath,amssymb}

\title{Note Title Goes Here}
\date{2026-01-01}

\begin{document}
\maketitle

Ordinary LaTeX prose goes here — your existing notes should mostly work
unchanged, including \verb|\newcommand| macros, \verb|itemize|/\verb|enumerate|,
\verb|figure| environments with \verb|\includegraphics|, and
\verb|verbatim|/code blocks. Math (inline \(...\) and display
\[...\] or \verb|equation|) is converted automatically and rendered with
KaTeX.

\section{A section heading}

For example:
\begin{equation}
H(q, p) = \frac{p^2}{2m} + V(q)
\end{equation}

\begin{figure}[h]
  \centering
  \includegraphics[width=0.6\textwidth]{../assets/img/your-figure.png}
  \caption{A caption explaining what the figure shows.}
\end{figure}

\end{document}

% ---------------------------------------------------------------------
% To publish this note:
%   1. Copy this file to a new name, e.g. my-new-note.tex (the filename
%      becomes the URL: notes/my-new-note.html).
%   2. Keep the "% tags:" and "% summary:" comment lines at the very top
%      — they're the only two pieces of metadata this build doesn't read
%      from standard LaTeX commands (\title, \date already work as-is).
%   3. Rebuild (requires pandoc — see README.md) — it will appear on the
%      Notes index automatically, sorted by date.
%
% Figures referenced with \includegraphics need to actually exist at that
% relative path (e.g. src/static/img/your-figure.png, referenced here as
% ../assets/img/your-figure.png since notes live one level under the
% published site's assets/ folder). LaTeX \label/\ref cross-references are
% not resolved by this pipeline — write figure/equation references as
% plain text instead.
%
% Files starting with an underscore (like this one) are never published —
% that's what keeps this safe to leave in place as a template.
% ---------------------------------------------------------------------
